\documentclass[12pt]{article}

\usepackage[utf8]{inputenc} % Pacote para permitir o uso de acentos
% Ajustes de hifenização
\tolerance=1
\emergencystretch=\maxdimen
\hyphenpenalty=10000
\hbadness=10000
\usepackage{lipsum} % Pacote para gerar texto aleatório
\usepackage{ragged2e} % Pacote para mais opções de alinhamento
\usepackage{graphicx} % Pacote para inserir figuras jpeg, pdf e png
\usepackage{ae}
\usepackage{epsfig}
\usepackage{mathtools}
\usepackage[export]{adjustbox}
\usepackage{lineno}
\usepackage{amsthm}
\usepackage{commath}
\usepackage{color, colortbl}
\usepackage{bbm}
\usepackage{setspace} % Pacote para ajustar espaçamento entre linhas

% Instructions to prevent LaTeX from hyphenating words
\tolerance=1
\emergencystretch=\maxdimen
\hyphenpenalty=10000
\hbadness=10000

\graphicspath{{figures/}} % Indica o diretório no qual as imagens estão sendo salvas, evitando que precisamos digitar parte do caminho da pasta que contém as imagens
\usepackage[rightcaption]{sidecap} % Pacote para mais opções de figura, colocar, título da legenda da figura ao lado direito
\usepackage{url} % Pacote para citar referências de sites
\usepackage{makeidx} % Pacote para criar índice
\makeindex % Habilitar a criação do índice
\usepackage{longtable} % Pacote para usar tabelas longas

% ----------
% Pacotes para dar quebra de linha dentro da tabela:
%\usepackage{fourier} 
%\usepackage{array}
\usepackage{makecell}

\renewcommand\theadalign{bc}
\renewcommand\theadfont{\bfseries}
\renewcommand\theadgape{\Gape[4pt]}
\renewcommand\cellgape{\Gape[4pt]}
% ----------

\usepackage{wrapfig} % Pacote para alinhar figuras à esquerda ou à direita do texto
\usepackage{ae}
\usepackage{epsfig}
\usepackage{mathtools}
\usepackage[export]{adjustbox}
\usepackage{lineno}
\usepackage{amsfonts}
\usepackage{amsthm}
\usepackage{commath}
\usepackage{color, colortbl}
\usepackage{bbm}

% ----------
% Pacote para inserir cabeçalho e rodapé:
% Atenção! Se remover o cabeçalho e rodapé, o LaTeX coloca como cabeçalho o padrão dele, que é bem interessante também!
% \usepackage{fancyhdr} % Pacote para inserir cabeçalho e rodapé
% \pagestyle{fancy} % Estilo do cabeçalho e rodapé, neste caso apenas uma linha no cabeçalho
%\lhead{{\footnotesize Prandiano\textregistered}} % Inserir texto no canto esquerdo do cabeçalho
%\chead{centro} % Inserir texto no centro do cabeçalho
%\rhead{{\footnotesize Curso de Python}} % Inserir texto no canto direito do cabeçalho
%\lfoot{{\footnotesize Copyright \textcopyright 2018 Prandiano}} % Inserir texto no canto esquerdo do rodapé
%\cfoot{centro} % Inserir texto no centro do rodapé
%\rfoot{{\footnotesize Reprodução proibida: Lei 9.610 de 19/02/1998}} % Inserir texto no canto direito do rodapé
% O comando \renewcommand re-nova o parâmetro especificado
%\renewcommand{\footrulewidth}{0.4pt} % Inserir uma linha acima do rodapé
% ----------

% ----------
% Pacote para inserir múltiplas figuras:
\usepackage{subfigure} % Pacote para inserir múltiplas figuras
% \usepackage{subcaption} % Pacote incompatível com subfigure
%\usepackage{caption}
% ----------

\usepackage{multicol} % Pacote que permite usar multicolunas
\usepackage{float} % Pacote para usar o "H" do figure
\usepackage[dvipsnames]{xcolor} % Pacote de cores para textos e possivelmente outros
\usepackage{amsmath} % Pacote para usar o cases
\usepackage[makeroom]{cancel} % Pacote para usar o cancelar (em e embaixo)
\usepackage{enumitem} % Pacote para enumerar itens
\usepackage{amssymb} % Pacote para inserir os símbolos dos conjuntos numéricos N, Z, Q, R, C
\usepackage[T1]{fontenc}
\usepackage{enumitem}
\usepackage{textcomp}
\usepackage{pdfpages} % Pacote para inserir página em PDF
\usepackage{courier} % Pacote para usar a fonte courier
%\usepackage{matlab-prettifier}
\usepackage{lscape} % Pacote para mudar sentido da página
\newtheorem{theorem}{Teorema} % Pacote para inserir Teorema
\newtheorem{corollary}{Corollary} % Pacote para inserir Corolário
\newtheorem{remark}{Remark}
\newtheorem{mytheorem}{Definition}
\usepackage{hyperref} % Pacote para inserir hyperref
\usepackage{multirow} % Pacote para usar os comandos \multirow
\usepackage{pdfpages} % Pacote para incluir PDF's
\usepackage{listingsutf8} % Pacote para inserir arquivos externos tipo .txt
\usepackage{etoolbox}
\AtBeginEnvironment{verbatim}{\footnotesize}

\usepackage{standalone} % Pacote para inserir sub-arquivos LaTeX
\usepackage{bookmark} % Pacote para NÃO tratar o Índice como Apêndice e colocar letra

% Nome da listagem em português
\renewcommand{\lstlistingname}{Listagem}

% REFERÊNCIAS: ----------
%\usepackage[numbers, compress]{natbib} % Pacote para usar as referências bibliográficas
%\usepackage{usebib}
% \bibliographystyle{references} % Estilo das referências
%\setcitestyle{authoryear,aysep={},open={(},close={)}}
%\newcommand{\myand}{e}

% ----------


% PARÂMETROS DA FUNÇÃO LISTINGS: ----------
\lstset
{
    basicstyle=\small\ttfamily,
    commentstyle=\color{Green},
    keywordstyle=\color{Cerulean},
    frame=single,
    language=Python,
    morekeywords={True, False},
    numbers=left,
    numbersep=10pt,
    numberstyle=\footnotesize\color{Gray},
    showstringspaces=false,
    stringstyle=\color{Mulberry},
    tabsize=3,
    breaklines=true, % Este comando é para quebrar uma linha comprida
    postbreak=\mbox{\textcolor{red}{$\hookrightarrow$}\space}, % Este comando é para indicar com uma seta vermelha que a linha continua embaixo 
    extendedchars=true,
    literate=
        {á}{{\'a}}1 {Á}{{\'A}}1 {à}{{\`a}}1 {À}{{\`A}}1 {ã}{{\~a}}1 {Ã}{{\~A}}1 {â}{{\^a}}1 {Â}{{\^A}}1
        {ç}{{\c{c}}}1 {Ç}{{\c{C}}}1 {é}{{\'e}}1 {É}{{\'E}}1 {ê}{{\^e}}1 {Ê}{{\^E}}1 {í}{{\'\i}}1 {Í}{{\'I}}1
        {ó}{{\'o}}1 {Ó}{{\'O}}1 {ô}{{\^o}}1 {Ô}{{\^O}}1 {õ}{{\~o}}1 {Õ}{{\~O}}1 {ú}{{\'u}}1 {Ú}{{\'U}}1
        {ü}{{\"u}}1 {Ü}{{\"U}}1 {“}{{``}}1 {”}{{''}}1 {’}{{'}}1 {—}{{--}}1 {–}{{--}}1 {−}{{-}}1
        {±}{{$\pm$}}1 {σ}{{$\sigma$}}1 {→}{{$\to$}}1 {…}{{\ldots}}1
}

\lstset{inputencoding=utf8/latin1}

\lstdefinelanguage{Matlab}{
	language=Matlab,
	basicstyle=\ttfamily\small,
	keywordstyle=\color{blue}\bfseries,
	stringstyle=\color{red},
	commentstyle=\color{gray}\itshape,
	morekeywords={function, end, if, else, for, while, switch, case, otherwise},
	sensitive=true
}
% ----------

% ----------
\usepackage{nomencl} % Pacote para criar nomenclatura e lista de nomenclatura
\makenomenclature % Criar a nomenclatura
\renewcommand{\nomname}{Lista de Nomenclaturas} % Renomear o título da nomenclatura
% ----------

\usepackage{amsmath} % Pacote para inserir matrix com chaves, entre outras coisas

\usepackage{booktabs} % Pacote para usar \toprule, \midrule e \bottomrule


% EXECUTAR CÁLCULOS A PARTIR DA LINHA DE COMANDO: ---------
%\usepackage{siunitx, xfp} % Pacote para executar cálculos a partir da linha de comando
%\usepackage{l3regex}%unless you have very new expl3
%%\sisetup{locale=DE}
%%\let\oldnum\num
%%\ExplSyntaxOn
%\def\num#1{\def\tmp{#1}\regex_replace_all:nnN{,([0-9])}{.\1}{\tmp}\oldnum{\tmp}}
%\ExplSyntaxOff
% ----------


%\pagecolor{black} % Alterar a cor da página
%\color{white} % Alterar a cor do texto

% ----------
% Pacotes para inserir algoritmos e pseudo-códigos:
\setlength{\marginparwidth}{2.5cm}
\usepackage[colorinlistoftodos]{todonotes}
\usepackage{algorithm} % Pacote para inserir algoritmos
\usepackage{algpseudocode} % Pacote para inserir algoritmos e pseudo-códigos
\floatname{algorithm}{Algoritmo} % Renomear a palavra Algorithm para Procedimento
\renewcommand{\algorithmicrequire}{\textbf{entrada:}} % Renomear a palavra Input para Entrada
\renewcommand{\algorithmicensure}{\textbf{saída:}} % Renomear a palavra Output para Saída
\renewcommand{\algorithmicwhile}{\textbf{enquanto}} % Renomear a palavra While para Enquanto
\renewcommand{\algorithmicif}{\textbf{se}} % Renomear a palavra If para Se
\renewcommand{\algorithmicend}{\textbf{fim}} % Renomear a palavra End Output para Fim do
\renewcommand{\algorithmicprocedure}{\textbf{procedimento}} % Renomear a palavra Procedure para Procedimento
\renewcommand{\algorithmicfor}{\textbf{para}} % Renomear a palavra For para Para
\renewcommand{\algorithmicelse}{\textbf{senão}} % Renomear a palavra Else para Senão
\renewcommand{\algorithmicfunction}{\textbf{função}} % Renomear a palavra Function para Função 
\renewcommand{\algorithmicthen}{\textbf{então}} % Renomear a palavra Then para Então
\renewcommand{\algorithmicdo}{\textbf{fazer}} % Renomear a palavra Do para Fazer
% ----------

\newcommand{\quotes}[1]{``#1''} % Comando para inserir um texto dentro de aspas-dupla com espaço após o final

% FLUXOGRAMA: ----------
\usepackage{tikz}
\usetikzlibrary{positioning, shapes.geometric}
\tikzstyle{line}=[draw] % here

\tikzstyle{startstop} = [rectangle, rounded corners, minimum width=3cm,
minimum height=1cm,text centered, draw=black, fill=white!30]
\tikzstyle{arrow} = [thick, ->, >=stealth]
\tikzstyle{process} = [rectangle, minimum width=3cm, minimum height=1cm, text centered, draw=black, fill=white!30]
\tikzstyle{io} = [trapezium, trapezium left angle=70, trapezium right angle=110, minimum width=3cm, minimum height=1cm, text centered, draw=black, fill=white!30]
\tikzstyle{connector} = [draw, circle, fill=white!20]
\tikzstyle{decision} = [diamond, minimum width=3cm, minimum height=1cm, text centered, draw=black, fill=white!30, aspect=3]
% ----------


% CAIXA ENTORNO DAS FÓRMULAS: ----------
\usepackage{dashbox}
\newcommand\dboxed[1]{\dbox{\ensuremath{#1}}}

% ----------

%\usepackage{tabularx} % para colunas com largura ajustável

% -------- CONFIGURAÇÃO DE MARGENS --------
% \usepackage[top=2cm,bottom=2cm,left=3cm,right=2cm]{geometry}

% -------- PACOTE fancyhdr --------
\usepackage{fancyhdr}

\setlength{\marginparwidth}{2cm}

% Ajuste do tamanho do cabeçalho para acompanhar exatamente \textwidth
\setlength{\headwidth}{\textwidth}

% Define um headheight mínimo (para evitar avisos)
\setlength{\headheight}{15pt} 

% Limpa configurações padrão de cabeçalho/rodapé
\pagestyle{fancy}
\fancyhf{} 

% Marcações automáticas de seção para o cabeçalho (opcional)
\renewcommand{\sectionmark}[1]{%
	\markboth{\thesection\ #1}{}%
}

% Texto do cabeçalho à esquerda: seção e título
% Usamos \leftmark pois o \sectionmark acima grava no \leftmark
\fancyhead[L]{\small \leftmark}

% Texto do cabeçalho à direita: número de página
\fancyhead[R]{\small \thepage}

% Linha horizontal no cabeçalho (espessura 0.4pt)
\renewcommand{\headrulewidth}{0.4pt}

% Sobrescreve o título padrão “References” para “Referências”
%\renewcommand{\refname}{Referências}
% Tradução de autoref para o português
\def\figureautorefname{Figura}
\def\tableautorefname{Tabela}
\def\equationautorefname{Equa\c{c}\~ao}

% \hypersetup{backref=true}

\ifdefined\backref
        \renewcommand*{\backref}[1]{}
        \renewcommand*{\backrefentrycount}[2]{#1\ifnum#2>1 ~(#2)\fi}
\fi

\DeclarePairedDelimiter{\ceil}{\lceil}{\rceil}

\onehalfspacing

\usepackage{pgfgantt}
  % OK
%++++++++++++++++++++++++++++++++++++++++++++++++++++++++++++++++++++++++++++++
% Identificacoes (se o trabalho for em inglês, insira os dados em inglês)
% Para entradas abreviadas de Professora (Profa.) em português escreva: Prof$^\textnormal{a}$.
%++++++++++++++++++++++++++++++++++++++++++++++++++++++++++++++++++++++++++++++
\course{Engenharia Aeronáutica e Mecânica} % Programa de PG ou Curso de Graduação
\area{Propulsão Aeroespacial e Energia} % Área de concentração na PG (Não utilizado no caso de TG)

% Autor do trabalho: Nome Sobrenome
\authorgender{masc} %sexo: masc ou fem
\author{\textbf{EDEN DENIS} Ferreira da Silva Lopes}{Santos}
\itaauthoraddress{Rua João Fonseca dos Santos, 20, Apto. 93}{12.230-088}{São José dos Campos--SP}

% Titulo da Tese/Dissertação
\title{Aplicação de Inteligência Artificial (AI) em Banco de Dados de Emissão de Motores}

% Orientador
\advisorgender{masc} % masc ou fem
\advisor{Prof.~Dr.}{Pedro Teixeira \textbf{LACAVA}}{ITA}

% Coorientador (Caso não haja coorientador, colocar ambas as variáveis \coadvisorgender e \coadvisor comentadas, com um % na frente)
\coadvisorgender{masc}
\coadvisor{}{}  % {Prof.~Dr.}{Carlos Alberto Rocha \textbf{PIMENTEL}}{ITA}

% Pró-reitor da Pós-graduação
\bossgender{fem}
\boss{Prof.~Dra.}{Emília \textbf{VILLANI}}

%Coordenador do curso no caso de TG
\bosscoursegender{masc}
\bosscourse{Prof.~Dr.}{ - } % masc ou fem

% Palavras-Chaves informadas pela Biblioteca -> utilizada na CIP
\kwcip{Engine taxonomy}
\kwcip{Emission classification}
\kwcip{Inteligência Artificial}
\kwcip{Artificial Intelligence}
\kwcip{Python}

% membros da banca examinadora
\examiner{Prof. Dr.}{Pedro Teixeira \textbf{LACAVA}}{President}{ITA}
\examiner{Prof. Dr.}{}{Internal Member}{ITA} % Vinícius \textbf{MALATESTA}}{Internal Member}{ITA}
\examiner{Prof. Dr.}{}{Internal Member}{ITA} % {Elisan dos Santos \textbf{MAGALHÃES}}{Internal Member}{ITA}
\examiner{Prof. Dr.}{}{Internal Member}{ITA} % {Carlos Alberto Rocha \textbf{PIMENTEL}}{External Member}{UFABC}
\examiner{Prof. Dr.}{}{Internal Member}{ITA} % {\textbf{KARL} Peter Burr}{External Member}{UFABC}

% Data da defesa (mês em maiúsculo, se trabalho em inglês, e minúsculo se trabalho em português) 
\date{14}{JULHO}{2025}

% Número CDU - (somente para TG)
\cdu{621.38}

% Folha de Registro do Documento
\newcommand{\advisorTitle}{Prof.~Dr.}
\newcommand{\advisorName}{Pedro Teixeira \textbf{LACAVA}}

\newcommand{\coadvisorTitle}{Prof.~Dr.}
\newcommand{\coadvisorName}{Carlos Alberto Rocha \textbf{PIMENTEL}}

\newcommand{\dateDay}{14}
\newcommand{\dateMonth}{JULHO}
\newcommand{\dateYear}{2025}

% Data de publicação e defesa (usada na folha de registro)
\newcommand{\frdDataPublicacao}{04 de agosto de 2025}
\newcommand{\frdNumeroRegistro}{DCTA/ITA/TD-029/2020}
\newcommand{\frdPalavrasAutor}{Big Data; Engine Emission; Emission Databank; Clustering; Artificial Intelligence}
\newcommand{\frdPalavrasResultado}{Taxonomia; Emissão de Motores; Planejamento de Dados; Inteligência Artificial; Processamento de Dados}
\newcommand{\frdApresentacao}{
	ITA, São José dos Campos. Curso de Mestrado. Programa de Pós-Graduação em Engenharia Aeronáutica e Mecânica. Área de Propulsão Aeroespacial e Energia.\\
	Orientador: \advisorTitle~\advisorName.\\
	Coorientador: \coadvisorTitle~\coadvisorName.\\
	Defesa em \dateDay/\dateMonth/\dateYear. Publicada em \frdDataPublicacao.
}

\newcommand{\PorcentagemDaLargura}{0.80}  % OK

\begin{document}


	\chapter{Introdução}

		Este capítulo apresenta o contexto da pesquisa, o problema de classificação de motores aeronáuticos e a motivação para aplicar métodos de agrupamento aos dados de emissões.
	
		\section{Introdução}
		
			% Eden:
			O problema do particionamento tem grande expressão também na área da Biologia, motivo pelo qual os cientistas criaram um ramo específico chamado de taxonomia, para essa área do conhecimento, a de organismos e plantas. Entretanto, tudo isso, teve início desde a criação das sociedades quando, Aristóteles por volta de $ 300 $ a. C. inicialmente separou, de forma generalista, os seres vivos das plantas e, a \textit{posteriori}, para o caso do primeiro grupo, por exemplo, separou-os pela cor do sangue entre vermelho e não. No século XVIII, Carl von Linné ($ 1707 $ - $ 1778 $), aprimorou o processo taxonômico (utilizado até hoje) que organizou melhor os seres vivos, pois já se conhecia uma quantidade maior de espécies. Atualmente  particionado em cinco reinos: \textit{Monera}, \textit{Protista}, \textit{Fungi}, \textit{Animalia} e \textit{Plantae}; Linné dividiu os seres vivos levando em conta sete critérios (reino, filo, classe, ordem, família, gênero e espécie), o primeiro critério mais generalista e, a medida que se avança o critério, faz-se uma análise mais específica, logo, tem-se uma melhor classificação. Já que o objetivo deste trabalho é classificar as aeronaves, tanto civis como militares, o neologismo Aerotaxonomia torna-se bastante adequado, uma vez que, a palavra aero está relacionada à Engenharia aeroespacial e, por conseguinte, taxonomia significa classificação \cite{Linnaeus1758}, tendo em vista que um termo que seria conveniente para a biologia seria biotaxonomia. Recentemente, foi adicionado um domínio antes do reino. Tendo em vista que a palavra aero (já relacionada) é abrangente, fica conveniente usar, para este caso, Aerotaxonomia de aeronaves (veículos) dado que, poderia ser também de motores (sistema propulsivo). 
			
			% Daqui para baixo quem escreveu foi o prof. Ricieri
			A Inteligência Artificial (AI), em particular suas duas técnicas ou métodos para  tratar dados, $ k $ -means e Dendrogramação, podem servir de modelo para o conhecimento moderno caracterizado pela sua limitada validade temporal. E uma Matemática como a do século XX, que contraria essa particularidade do mundo moderno, século XXI, sociedade liquida de Bauman \cite{Bauman2000}, é inútil para os propósitos de toda e qualquer instituição do saber científico-tecnológico atual e também, muito provável, do futuro.
			
			Para o cientista de dados, a Inteligência Artificial (AI) e suas técnicas é o que deve ser estudado e aplicado na Engenharia moderna, em particular na aeronáutica; não porque ofereça um campo novo para as investigações do conhecimento, como aconteceu com a Matemática clássica até então, mas sim porque seus resultados preditivos encorajam a elevar o \textit{status} da Inteligência Artificial (AI) como a principal disciplina que poderá governar as Engenharias, essas que deverão, obrigatoriamente, preocuparem-se com aspectos nunca antes pensados pelos engenheiros como preferências populares, planos de venda, opções de \textit{marketing}, percentuais de sucesso e subjetividades de formas do produto; isso tudo exigido também pelos produtos aeronáuticos.  
			
			%		Diante dessas exigências do mercado em relação a profissão de Engenharia não existe outro caminho se não o da Inteligência Artificial (AI), que possa ajustar os currículos universitários às demandas atuais. 
			
			Enquanto o engenheiro tradicional busca a todo custo impor aos fenômenos do mundo teorias que possam explicá-los completamente; o engenheiro de dados analisa-os como eles se apresentam a nós, com certo grau de ambiguidade. Essa abordagem nova, pragmática, que tolera possíveis erros ou desvios de previsão, do ponto de vista semiótico–abducional, alegoriza dois movimentos de predições, de mesma direção e sentidos inversos: dos modelos, com seus símbolos sintáxicos aos dados e com suas minerações computacionais aos modelos. Daí a essência da dissonância entre Engenharia clássica, primeiro sentido e de Engenharia de dados, segundo sentido, estar nas explicações teóricas completas e nos modelos computacionais limitados. 
			
			Para um engenheiro de dados, profissional de natureza imediatista, já que atua em uma sociedade cada vez mais liquefeita, a busca de um modelo perfeito, oficial e universalmente abrangente, proíbe-o de qualquer possibilidade de predizer com liberdade, possíveis resultados minimamente assertivos. Assim, a essência do conflito entre esses dois profissionais das exatas deslocados no tempo, resvala na diversidade do conhecimento produzido nos inúmeros setores dessa ou daquela Engenharia, que exigem respostas rápidas e objetivas, mas não extremamente precisas. 
			
			Desta forma, sabendo que a Inteligência Artificial (AI) enquanto Matemática que tem influenciado como nunca antes as Engenharias, permitindo inúmeras respostas, resta àquele que se diz engenheiro ou cientista de dados saber quais são as perguntas. 
			
			% Daqui para baixo quem escreveu foi o Eden
			Portanto, a partir desta Matemática moderna, ajustada a sociedade atual e que tem definido novos padrões de solução somado à taxonomia, o principal objetivo deste trabalho é descrever um processo para particionar uma população $ d $-dimensional de aeronaves em $ K $-sets, sejam estas civis ou militares. A modelagem Matemática utilizada é chamada de $ k $ -means que, trata-se de uma aproximação eficiente quando há variações de determinadas características entre as classes, porém, reduzindo uma possível subjetividade humana. Apesar do interesse ser minimizar possíveis subjetividades na classificação das aeronaves, deve-se corroborar a análise Matemática com a prática da experiência computacional. O método escolhido para particionar, $ k $ -means, é um processo de fácil desenvolvimento e computacionalmente econômico, portanto, factível sobretudo pelo fato de não exigir investimento de altíssimos custos. 
			
			Além de toda a aplicação envolvida, o método $ k $ -means que, originalmente foi concebido como uma tentativa de encontrar um procedimento computacional factível para classificação otimizada de problemas da biologia e propriamente da computação \cite{MacQueen1967}, também é teoricamente pertinente, o que levou a escolhê-lo como ferramenta Matemática deste trabalho. Este método é uma generalização para a média ordinária da amostra, logo, é possível analisar o comportamento assintótico pertinente as características das classes e determinar algum tipo de lei baseada na Inteligência Artificial (AI) das aeronaves quando analisa-se o $ K $ do $ k $ -means para um número grande. O modelo matemático $ k $ -means é desenvolvido na seção (\ref{subsec:kmeans}) e os resultados obtidos sobre o comportamento assintótico são descritos na seção (\ref{sec:ResultadosEDiscussoes}).
			
			Em geral, o método não converge para uma solução ótima global, embora possivelmente existam casos especiais onde isso acontece. Entretanto, até onde se sabe, não existe método factível que sempre se obtêm solução global. Sendo assim, garante-se que pelo menos um resultado falibalista será encontrado, solução ótima local \cite{MacQueen1967}.
	
\end{document}
